\section{USAMO 2015}
        \begin{enumerate}
        	\item (\textit{USAMO 2015}) Solve in integers the equation

 \[x^2+xy+y^2 = \left(\frac{x+y}{3}+1\right)^3.\]


 

	\item (\textit{USAMO 2015}) Quadrilateral $APBQ$ is inscribed in circle $\omega$ with $\angle P = \angle Q = 90^{\circ}$ and $AP = AQ < BP$. Let $X$ be a variable point on segment $\overline{PQ}$. Line $AX$ meets $\omega$ again at $S$ (other than $A$). Point $T$ lies on arc $AQB$ of $\omega$ such that $\overline{XT}$ is perpendicular to $\overline{AX}$. Let $M$ denote the midpoint of chord $\overline{ST}$. As $X$ varies on segment $\overline{PQ}$, show that $M$ moves along a circle.


 

	\item (\textit{USAMO 2015}) Let $S = \{1, 2, ..., n\}$, where $n \ge 1$. Each of the $2^n$ subsets of $S$ is to be colored red or blue. (The subset itself is assigned a color and not its individual elements.) For any set $T \subseteq S$, we then write $f(T)$ for the number of subsets of T that are blue.


 Determine the number of colorings that satisfy the following condition: for any subsets $T_1$ and $T_2$ of $S$,

 \[f(T_1)f(T_2) = f(T_1 \cup T_2)f(T_1 \cap T_2).\]


 

	\item (\textit{USAMO 2015}) Steve is piling $m\geq 1$ indistinguishable stones on the squares of an $n\times n$ grid. Each square can have an arbitrarily high pile of stones. After he finished piling his stones in some manner, he can then perform stone moves, defined as follows. Consider any four grid squares, which are corners of a rectangle, i.e. in positions $(i, k), (i, l), (j, k), (j, l)$ for some $1\leq i, j, k, l\leq n$, such that $i<j$ and $k<l$. A stone move consists of either removing one stone from each of $(i, k)$ and $(j, l)$ and moving them to $(i, l)$ and $(j, k)$ respectively, or removing one stone from each of $(i, l)$ and $(j, k)$ and moving them to $(i, k)$ and $(j, l)$ respectively.


 Two ways of piling the stones are equivalent if they can be obtained from one another by a sequence of stone moves.


 How many different non-equivalent ways can Steve pile the stones on the grid?


 

	\item (\textit{USAMO 2015}) Let $a, b, c, d, e$ be distinct positive integers such that $a^4 + b^4 = c^4 + d^4 = e^5$. Show that $ac + bd$ is a composite number.


 

	\item (\textit{USAMO 2015}) Consider $0<\lambda<1$, and let $A$ be a multiset of positive integers. Let $A_n=\{a\in A: a\leq n\}$. Assume that for every $n\in\mathbb{N}$, the set $A_n$ contains at most $n\lambda$ numbers. Show that there are infinitely many $n\in\mathbb{N}$ for which the sum of the elements in $A_n$ is at most $\frac{n(n+1)}{2}\lambda$. (A multiset is a set-like collection of elements in which order is ignored, but repetition of elements is allowed and multiplicity of elements is significant. For example, multisets $\{1, 2, 3\}$ and $\{2, 1, 3\}$ are equivalent, but $\{1, 1, 2, 3\}$ and $\{1, 2, 3\}$ differ.)


 

\end{enumerate}